\section{Introdução}
\label{sec:introducao}

Linhas de produto de software representam famílias de produtos sobre uma base comum, e feature models são o padrão \emph{de facto} para descrever sua variabilidade~\cite{galindo2023llm}. Modelos industriais podem ultrapassar 7.000 features (p.ex., o kernel do Linux), o que torna a análise automatizada indispensável. Avaliar essas ferramentas, contudo, exige grandes conjuntos de modelos realistas, que raramente são públicos por serem estratégicos para as empresas. Geradores existentes produzem modelos sintaticamente válidos, mas com nomes sem significado (F1, F2, \ldots).

Galindo et al.~\cite{galindo2023llm} propuseram usar um LLM para estender modelos reais escritos em \uvl{}~\cite{sundermann2021uvl}, gerando variantes que preservem sintaxe e domínio. O estudo usou o Codex~\cite{chen2021codex}, desativado pela OpenAI em 2023, e publicou as respostas obtidas.

Este trabalho apresenta (i)~uma \textbf{reprodução} da análise com os artefatos originais e (ii)~uma \textbf{replicação} com LLMs atuais, segundo a classificação de Gómez et al.~\cite{gomez2014replication}. \todo{contribuições em itens; organização do artigo.}
